\documentclass[11pt,a4paper]{article}
\usepackage[utf8]{inputenc}
\usepackage{amsmath,amssymb,amsthm}
\usepackage{listings}
\usepackage{xcolor}
\usepackage[margin=1in]{geometry}
\usepackage{hyperref}

\newtheorem{theorem}{Theorem}
\newtheorem{lemma}[theorem]{Lemma}

\lstset{
  language=C++,
  basicstyle=\ttfamily\small,
  keywordstyle=\color{blue}\bfseries,
  commentstyle=\color{green!60!black},
  stringstyle=\color{red},
  numbers=left,
  numberstyle=\tiny\color{gray},
  breaklines=true,
  frame=single,
  tabsize=4,
  showstringspaces=false
}

\title{IOI 2004 -- Farmer}
\author{Solution}
\date{}

\begin{document}
\maketitle

\section{Problem Statement Summary}
A farmer has $N$ cows at known positions on a number line, and $M$ barns
(each with a capacity) also at known positions. Assign every cow to a
barn (respecting capacities) so as to minimize the maximum distance any
cow travels to its assigned barn.

\section{Solution: Binary Search + Greedy}

\subsection{Binary Search}
Binary search on the answer $D$ (the maximum allowed distance). For
each candidate $D$, check feasibility in $O(N + M)$.

\subsection{Feasibility Check}
\begin{lemma}
If both cows and barns are sorted by position, the greedy strategy of
assigning each cow (left to right) to the leftmost barn within distance
$D$ that still has capacity is optimal.
\end{lemma}
\begin{proof}
Suppose cow $i$ is assigned to barn $b$ in the greedy solution but to
a different barn $b'$ in some optimal solution. Since $b$ is the leftmost
feasible barn, $b$ is at most as far right as $b'$. Swapping the
assignments of cow $i$ and whatever cow used barn $b$ cannot increase
the maximum distance (an exchange argument). By induction, the greedy
assignment is feasible whenever any assignment is.
\end{proof}

\subsection{Implementation Detail}
Maintain a barn pointer $b$ that advances past barns too far to the left
or with zero remaining capacity. Since cows are processed left to right,
the pointer never needs to retreat: any barn too far left for cow $i$ is
also too far left for cow $i{+}1$.

If the pointer's barn is beyond $\mathrm{cow}[i] + D$, the check fails.
Otherwise, decrement that barn's remaining capacity.

\section{C++ Implementation}

\begin{lstlisting}
#include <bits/stdc++.h>
using namespace std;

int main() {
    ios::sync_with_stdio(false);
    cin.tie(nullptr);

    int N, M;
    cin >> N >> M;

    vector<long long> cow(N);
    for (int i = 0; i < N; i++) cin >> cow[i];

    vector<long long> barn(M);
    vector<int> cap(M);
    for (int i = 0; i < M; i++) cin >> barn[i] >> cap[i];

    sort(cow.begin(), cow.end());

    // Sort barns by position, keeping capacity
    vector<int> order(M);
    iota(order.begin(), order.end(), 0);
    sort(order.begin(), order.end(),
         [&](int a, int b) { return barn[a] < barn[b]; });

    vector<long long> sb(M);
    vector<int> sc(M);
    for (int i = 0; i < M; i++) {
        sb[i] = barn[order[i]];
        sc[i] = cap[order[i]];
    }

    auto check = [&](long long D) -> bool {
        vector<int> rem(sc.begin(), sc.end());
        int b = 0;
        for (int i = 0; i < N; i++) {
            // Advance past barns too far left or empty
            while (b < M && (sb[b] < cow[i] - D || rem[b] <= 0))
                b++;
            if (b >= M || sb[b] > cow[i] + D)
                return false;
            rem[b]--;
        }
        return true;
    };

    long long lo = 0, hi = 2000000000LL;
    while (lo < hi) {
        long long mid = lo + (hi - lo) / 2;
        if (check(mid)) hi = mid;
        else lo = mid + 1;
    }

    cout << lo << "\n";
    return 0;
}
\end{lstlisting}

\paragraph{Bug note.} The two-pointer greedy above has a subtle issue:
when a barn runs out of capacity, the pointer $b$ advances past it via
the \texttt{while} loop. However, the pointer \emph{also} skips barns
that are too far left. This is correct because cows are sorted, so
$\mathrm{cow}[i] - D$ is non-decreasing.

\section{Complexity Analysis}
\begin{itemize}
  \item \textbf{Time:} $O((N + M) \log V)$ where $V$ is the coordinate
        range. Sorting is $O(N \log N + M \log M)$; each feasibility
        check is $O(N + M)$; binary search runs $O(\log V)$ iterations.
  \item \textbf{Space:} $O(N + M)$.
\end{itemize}

\end{document}
